\chapter{补充实验与附录材料}

\section{候选对比方向与扩展比较}\label{app:candidate-schemes}

\subsection{与 JIT 时机和解释器路径相关的候选方向}

本节列出与 JIT 时机和解释器路径相关、但未作为正文正式机制展开的候选方向。这些方向大多具有明确的局部热点依据和实现对象，但在当前实验设置下，尚未形成稳定且可推广的 benchmark 级性能提升。因此，本节保留这些方向，主要用于补充说明正式机制的选择依据，并展示不同技术路线的适用边界。

\begin{figure}[htbp]
    \centering
    \includegraphics[width=0.88\textwidth]{mybenchmark/candidate-schemes-heatmap.pdf}
    \caption{附录对比方向在固定对照版本下的归一化比较}
    \label{fig:candidate-schemes-heatmap}
\end{figure}

图~\ref{fig:candidate-schemes-heatmap} 展示了固定对照版本下多个候选方向的比较结果。图中 \texttt{1.0} 表示基准性能，低于该值表明该方向未形成性能提升。候选方向涵盖阈值调节、OSR 边界改写、解释器快路径以及通用复制入口优化等技术路线\citep{Arnold2005,Lameed2013,Haubl2008OptimizedStrings,Lemire2022UnicodeSIMD}。这些结果说明，局部热点的可见性并不必然转化为系统级性能提升。

\begin{table}[htbp]
    \centering
    \small
    \caption{JIT 时机与解释器路径相关候选方向的比较结果}
    \label{tab:appendix-jit-candidates}
    \begin{tabular}{>{\raggedright\arraybackslash}p{3.8cm}>{\raggedright\arraybackslash}p{2.2cm}>{\raggedright\arraybackslash}p{5.0cm}}
        \hline
        \textbf{候选方向} & \textbf{典型结果} & \textbf{适用性与局限性} \\
        \hline
        \texttt{optimize\_threshold} 调整 & \texttt{pmd -11.23\%} & 作用于第二层 optimizing recompilation，介入时机偏后，难以改善短生命周期负载的首层编译瓶颈 \\
        JIT OSR 热路径内联调整 & \texttt{pmd -17.78\%} & 更贴近解释器/OSR 局部路径，无法缓解当前主导瓶颈 \\
        Nterp 热点循环 OSR 重检查调整 & \texttt{mpegaudio -1.84\%} & 引入额外固定开销，缺乏稳定性能提升 \\
        解释器已解析类型快路径 & \texttt{pmd -18.39\%} & 改善局部解析效率，但附加判断增加了主路径成本 \\
        JIT debug-info 延后生成 & 多数工作负载低于对照版本 & 存在局部热点，但难以建立微观生效机制与宏观 benchmark 性能提升之间的稳定联系 \\
        \hline
    \end{tabular}
\end{table}

如表~\ref{tab:appendix-jit-candidates} 所示，这类候选方向的共同问题在于：要么作用层级偏后，要么只改善局部执行路径，要么通过额外检查和代码膨胀引入新的固定成本。因此，在当前实验条件下，它们尚不足以替代正文采用的基线 JIT 提早触发机制。

\subsection{与运行时热路径相关的候选方向}

\begin{table}[htbp]
    \centering
    \small
    \caption{运行时热路径相关候选方向的比较结果}
    \label{tab:appendix-runtime-candidates}
    \begin{tabular}{>{\raggedright\arraybackslash}p{4.0cm}>{\raggedright\arraybackslash}p{2.2cm}>{\raggedright\arraybackslash}p{4.8cm}}
        \hline
        \textbf{候选方向} & \textbf{典型结果} & \textbf{适用性与局限性} \\
        \hline
        原始数组批量搬运展开 & \texttt{pmd -9.79\%} & 泛化复制操作并非目标工作负载的主要开销 \\
        \texttt{memcpy} 快速入口接线 & \texttt{pmd -11.95\%} & 目标负载较少稳定触发该底层入口 \\
        CRC32 JNI 快路径 & \texttt{mpegaudio -14.28\%} & 热点占比过低，缺乏全局优化意义 \\
        \texttt{DexFileVerifier} 连续栈容器替换 & \texttt{pmd -11.91\%} & 单点替换未能覆盖启动期完整执行链，已在正文中改写为阶段链优化问题 \\
        命令行参数/JIT init 相关方向 & 多数结果低于对照版本 & 仅降低局部解析开销，未触及主导性能瓶颈 \\
        \hline
    \end{tabular}
\end{table}

上述结果进一步表明，运行时热路径中的局部热点并不一定具有全局优化价值。以 CRC32 JNI 快路径为例，尽管其函数入口明确，但在目标工作负载中的开销占比过低。类似地，\texttt{memcpy} 优化与原始数组批量搬运虽然涉及复制操作，却未像正文中的运行时热路径内存批量搬移机制那样，稳定覆盖字符串替换、拼接与 UTF-16 构造等核心执行路径。

针对 \texttt{DexFileVerifier} 的单一容器替换也未能形成独立性能提升，因此本文未将其保留为单点修补方案，而是在正文中转而采用沿 container probing、UTF-8 校验与 LEB128 解码主链展开的阶段链式优化。

\subsection{候选方向的共性边界}

结合图~\ref{fig:candidate-schemes-heatmap}、表~\ref{tab:appendix-jit-candidates} 与表~\ref{tab:appendix-runtime-candidates}，这些候选方向主要受限于三类因素。

第一，部分方向仅缓解局部热点，但未覆盖主导执行路径。此类优化在 flamegraph 中通常可见，却难以转化为宏观 benchmark 层面的稳定性能提升。第二，部分方向的生效时机偏后或性能提升窗口与目标工作负载不匹配。例如，第二层编译时机调整与 OSR 边界改写，并不能解决短生命周期负载最敏感的首层编译介入问题。第三，局部开销的降低可能破坏系统整体平衡。某些方向虽然减少了特定计算，却引入了冗余判断、增加了代码体积或降低了路径复用率，最终拖累整体性能。

这些扩展比较从侧面说明，正文保留的正式机制并非由局部热点直接决定，而是建立在主路径命中、机制可解释性与稳定性能提升三者同时成立的基础上。


\section{Intrinsic 机制的补充验证}\label{app:intrinsic-ablation}

\subsection{实验设计与比较设置}

本节补充分析 LoongArch64 架构下 ART intrinsic 机制的整体性能提升边界。由于 intrinsic 优化主要涉及广泛的库调用与编译器底层机制扩展，超出了正文聚焦的几类核心问题，故未将其列为正式机制。本节将其作为补充比较，用以说明 LoongArch64 平台上 JIT 与 optimizing 路径的性能改善并不只来源于正文中的几类机制。

实验分别在 \texttt{JIT} 与 \texttt{interpreter} 模式下评估开启与关闭 intrinsic 的性能差异。其中，\texttt{JIT} 模式用于观察 intrinsic 的实际性能提升，而 \texttt{interpreter} 模式作为对照组。若性能增益主要来自编译器的 intrinsic 分发与替换机制，则 \texttt{interpreter} 模式下的性能比值应接近 \texttt{1.0x}。

\subsection{几何平均结果}

\begin{figure}[htbp]
    \centering
    \includegraphics[width=0.78\textwidth]{mybenchmark/intrinsic-ablation-geomean.pdf}
    \caption{Intrinsic 开关在 \texttt{JIT} 与 \texttt{interpreter} 下的几何平均结果}
    \label{fig:intrinsic-ablation-geomean}
\end{figure}

如图~\ref{fig:intrinsic-ablation-geomean} 所示，在 \texttt{JIT} 模式下，LoongArch64 intrinsic 机制带来了稳定的整体性能提升。根据式~\ref{eq:geomean-speedup} 计算，所有有效工作负载的几何平均加速比为 \texttt{1.0415x}。其中，\texttt{SPEC startup}、\texttt{SPEC throughput} 与 \texttt{DaCapo} 的加速比分别为 \texttt{1.0339x}、\texttt{1.0508x} 与 \texttt{1.0387x}。作为对照，\texttt{interpreter} 模式的结果大多维持在 \texttt{\pm 2\%} 的波动范围内。

这些结果表明，intrinsic 的性能增益主要来自 \texttt{JIT/optimizing} 编译路径，而非环境波动或其他构建差异。虽然 intrinsic 并非本文的正式研究对象，但这一补充结果有助于完善 LoongArch64 平台 ART 的整体性能画像。

\subsection{逐工作负载边界分析}

\begin{figure}[htbp]
    \centering
    \includegraphics[width=0.88\textwidth]{mybenchmark/intrinsic-ablation-heatmap.pdf}
    \caption{Intrinsic 开关在代表性 workload 上的逐项比较}
    \label{fig:intrinsic-ablation-heatmap}
\end{figure}

进一步由图~\ref{fig:intrinsic-ablation-heatmap} 可见，intrinsic 的优化效果高度集中于加解密任务及部分依赖密集库调用的工作负载，而通用控制流密集型负载则几乎没有性能提升。具体而言，\texttt{crypto.aes} 与 \texttt{startup.crypto.aes} 的加速比分别达到 \texttt{1.4996x} 和 \texttt{1.2993x}，\texttt{h2} 与 \texttt{crypto.signverify} 也分别获得了 \texttt{1.1513x} 和 \texttt{1.0604x} 的提升。相比之下，\texttt{pmd} 出现了轻微回退（\texttt{0.9474x}），说明该负载的主导成本并不位于 intrinsic 所覆盖的路径上。

总体而言，LoongArch64 平台 ART 的性能优化呈现出明显的分层结构。正文主要关注控制流、运行时热路径、启动期校验链和后端向量化等问题，而本节补充的 intrinsic 结果则展示了更通用的底层库优化潜力。


\section{完整 SIMD 覆盖候选版的扩展比较}\label{app:full-simd-ablation}

\subsection{S0、S1、S2 实验设计}

本节分析完整 SIMD 覆盖候选版本（\texttt{Snapshot-E}）的扩展比较结果。该对比旨在说明：为什么正文采用最小可用版本 \texttt{S1} 作为正式方案，同时将完整覆盖候选版本 \texttt{S2} 保留为扩展比较对象。

实验设定了三种对比状态：\texttt{S0} 代表无 SIMD 优化的基准版本；\texttt{S1} 为正文采用的最小可用 LSX 浮点 SIMD 路径；\texttt{S2} 为完整 SIMD 覆盖候选版本，补充了更多 lowering 规则与惯用语（idiom）识别支持。

这一比较的目的，不是重新定义正文正式方案，而是观察在扩大后端向量化覆盖范围后，性能提升是否能够从少数核心负载扩展到更广泛的工作负载集合。

\subsection{全部工作负载与关键工作负载结果}

\begin{figure}[htbp]
    \centering
    \includegraphics[width=0.88\textwidth]{mybenchmark/loongarch64-lsx-full-simd-ablation-geomean.pdf}
    \caption{\texttt{S0/S1/S2} 在不同工作负载集合上的几何平均比较}
    \label{fig:app-lsx-full-simd-ablation-geomean}
\end{figure}

\begin{figure}[htbp]
    \centering
    \includegraphics[width=0.88\textwidth]{mybenchmark/loongarch64-lsx-full-simd-ablation-key-speedup.pdf}
    \caption{\texttt{S0/S1/S2} 在关键工作负载上的相对性能变化比较}
    \label{fig:app-lsx-full-simd-ablation-key}
\end{figure}

图~\ref{fig:app-lsx-full-simd-ablation-geomean} 统计了各版本在全部 \texttt{JIT} 工作负载下的几何平均加速比（均基于式~\ref{eq:geomean-speedup} 进行同向换算）。整体而言，\texttt{S1} 的平均性能仍优于 \texttt{S2}：\texttt{S1/S0} 达到 \texttt{1.0318x}，\texttt{S2/S0} 为 \texttt{1.0252x}，而 \texttt{S2/S1} 降至 \texttt{0.9936x}。

不过，图~\ref{fig:app-lsx-full-simd-ablation-key} 的细粒度比较同时显示，\texttt{S2} 在 \texttt{lu.small}、\texttt{h2} 和 \texttt{avrora} 等关键工作负载上的表现反超了 \texttt{S1}。这说明更广泛的后端指令覆盖具备客观的优化潜力，只是其性能影响在不同负载之间呈现出更明显的分化。

基于这些结果，\texttt{S2} 更适合被视为一种能够在特定核心负载上进一步扩大性能改善幅度的候选增强版，而尚不足以在整体几何平均意义下替代 \texttt{S1}。

\subsection{正式方案与候选增强版的边界}

从边界分析的角度看，\texttt{S1} 与 \texttt{S2} 的差异主要体现在三个方面。

首先，二者的目标不同。\texttt{S1} 旨在验证最小可用 LSX 浮点路径能否形成稳定、可解释且可复现的性能提升；\texttt{S2} 则进一步探索在扩大指令覆盖与 idiom 识别范围后，是否能够获得额外性能空间。

其次，二者的适用范围不同。\texttt{S1} 在整体几何平均上更稳健，因而适合作为正文中的正式方案；\texttt{S2} 虽在若干关键负载上具有更强表现，但其性能改善在更广泛工作负载集合中的一致性仍不足。

最后，二者的研究定位不同。\texttt{S1} 用于支撑正文中关于最小可用 LSX 路径的正式结论；\texttt{S2} 则用于揭示在扩展向量化覆盖范围后，性能改善与稳定性之间的进一步权衡关系。

因此，正文采用 \texttt{S1} 作为正式方案，而将 \texttt{S2} 保留为附录中的扩展比较对象。这样的处理方式既能保持正文结论的稳定性，也能完整呈现更广泛 SIMD 覆盖的潜在性能提升与边界。
